The goal of this Concept of Operations is to provide a high-level description about what \abbr{EWAS} should achieve, to what extent, how it would be operated and what phases does the mission's life cycle consist of. 

\subsection{Mission objectives and requirements}

The primary purpose of the Early Warning Alarm System is to identify missile-launch events that could threaten the Baltic region, including launches originating outside it, derive impact-area estimates from satellite observations, and send secure warning products to authorized national authorities. 

\textbf{Main objectives}

\begin{itemize}
    \item Detecting ballistic missile events threatening the Baltic region at the earliest possible stage of flight using infrared signatures. 

    \item Tracking and prediction of detected missile's trajectory to efficiently estimate the possible impact area. 

    \item Generating and dissemination clear warning information to authorities of potentially endangered areas enabling alerting civilians ahead of the impact. 

    \item Storing historical event data to support analysis and improve the system's prediction capability. 
\end{itemize}

\textbf{Main requirements}

\begin{itemize}
    \item Observation coverage adequate to identify relevant threats to the Baltic region, including launch events occurring outside that region; global coverage remains subject to supporting simulation evidence. 
    
    \item Low end-to-end latency between the space and ground segments ensuring sufficient time for users to react to the incoming missile threat before impact. 

    \item Accurate estimation of the preliminary impact area of the tracked missile. 

    \item High classification accuracy resulting in maximizing the missile detection probability and minimizing the false alarm rate. 

    \item High availability, preventing time windows in which the system is "blind" and ensuring proper alert dissemination. 
\end{itemize}

\subsection{System characteristics}

From the perspective of the strategic operator, \abbr{EWAS} presents a concise, continuously updated view of the event. The Mission Control Center displays the candidate launch location, the refined trajectory estimate, the probabilistic impact area, and the confidence in those results. It uses clear visual and audible notifications. After the Ground Segment has received observations from the two selected \abbr{LEO} satellites and produced an impact assessment, the Early Warning System sends the first preliminary warning automatically to authorized authorities, without waiting for operator approval. The operator then reviews the event and can confirm, update, or cancel the preliminary warning through a controlled and audited workflow. Public alerting beyond the authorized authority interface is outside the \abbr{EWAS} boundary. 

\subsection{Mission operation}

The Early Warning Alarm System comprises only the Space Segment and the Ground Segment. It operates a continuous surveillance and event-response loop; governmental authority systems are external warning recipients, not a third system segment. 

\textbf{Space segment operation}

The Space Segment combines \abbr{GEO} detection and \abbr{LEO} tracking layers. Following a candidate thermal-event detection and onboard filtering, the detecting \abbr{GEO} satellite sends an event report to the Ground Segment. The Ground Segment selects two suitable \abbr{LEO} tracking satellites on the basis of the reported position, viewing geometry, and link availability, and sends them tracking commands through an available ground-station link or an \abbr{ISL} relay route where available. Inter-satellite communication remains possible when needed. The selected \abbr{LEO} satellites obtain complementary observations, form a preliminary trajectory estimate, and transmit the observations and estimate to the Ground Segment. The timeliness of this architecture is a verification item, not an assumed proven capability.

\textbf{Ground segment operation}

The Ground Segment receives and validates the initial \abbr{GEO} report, selects and cues two \abbr{LEO} satellites, and receives their encrypted observations. Ground Station Network functions handle reception and secure forwarding; the Data Processing System conducts the principal fusion and computes a refined trajectory and probabilistic impact-area estimate. Once these \abbr{LEO}-based results are available, the Ground Segment Early Warning System automatically sends a preliminary warning to the appropriate authorized national authority interface. The Mission Control Center displays the information for subsequent operator review and controlled follow-up notifications. The Data Archiving System maintains the required records. The end-to-end delivery target remains 60 seconds from initial \abbr{GEO} detection; meeting this target depends on command routing, tracking availability, downlink and processing performance and has not yet been demonstrated.

\subsection{Life cycle concept}

\begin{enumerate}
    \item \textbf{Concept definition}: The initial phase of the mission's life cycle consists of defining sensor parameters, designing the orbital constellation as well as preliminary testing and software validation. 

    \item \textbf{Production and deployment}: Satellite components are built in cleanroom facilities, subjected to rigorous testing and integrated together. Afterwards, the satellites are launched into their designated orbits and the ground segment is established. 

    \item \textbf{Operation and maintenance}: Following the checkout and calibration the satellites enter their active service life during which trained operators monitor the state of the system and react accordingly to any detected missile threat.  
    
    During this phase software upgrades are delivered to both the space and ground segments, while hardware maintenance is limited to the ground segment only. 

    \item \textbf{End-of-life and disposal}: At the end of operational life, to prevent space debris accumulation, \abbr{LEO} satellites are de-orbited and are left to burn-up safely in the Earth's atmosphere while \abbr{GEO} satellites are moved into graveyard orbit. 
\end{enumerate}
